% Part: normal-modal-logic
% Chapter: completeness
% Section: completeness-K

\documentclass[../../../include/open-logic-section]{subfiles}

\begin{document}

\olfileid[id]{nml}{com}{cmk}

\olsection{Penentuan dan Kelengkapan bagi \Log{K}}

Sekarang kita siap menggunakan model kanonik untuk menetapkan kelengkapan.
Kelengkapan mengikuti dari fakta bahwa !!{formula}s yang benar dalam model
kanonik bagi~$\Sigma$ tepat merupakan formula yang !!{derivable}
berdasarkan~$\Sigma$. Model yang mempunyai sifat ini dikatakan
\emph{menentukan}~$\Sigma$.

\begin{defn}
  Suatu model $\mModel{M}$ \emph{menentukan} suatu logika modal
  normal~$\Sigma$ tepat ketika $\mSat{M}{!A}$ jika dan hanya jika $\Sigma
  \Proves !A$, bagi semua !!{formula}s~$!A$.
\end{defn}

\begin{thm}[Penentuan]\ollabel{thm:determination}
   $\mSat{M^\Sigma}{!A}$ jika dan hanya jika $\Sigma \Proves !A$.
\end{thm}

\begin{proof}
  Jika $\mSat{M^\Sigma}{!A}$, maka bagi setiap himpunan lengkap dan
  $\Sigma$-konsisten~$\Delta$, berlaku
  $\mSat{M^\Sigma}{!A}[\Delta]$. Jadi, berdasarkan Lema Kebenaran,
  $!A \in \Delta$ bagi setiap himpunan lengkap dan
  $\Sigma$-konsisten~$\Delta$; karena itu, berdasarkan
  \olref[lin]{cor:provability-characterization} (dengan $\Gamma =
  \emptyset$), $\Sigma \Proves !A$.

  Sebaliknya, jika $\Sigma \Proves !A$, maka berdasarkan
  \olref[ccs]{prop:ccs-properties}\olref[ccs]{prop:ccs-closed}, setiap
  himpunan lengkap dan $\Sigma$-konsisten~$\Delta$ memuat~$!A$; sehingga,
  berdasarkan Lema Kebenaran, $\mSat{M^\Sigma}{!A}[\Delta]$ bagi setiap
  $\Delta \in W^\Sigma$, yakni $\mSat{M^\Sigma}{!A}$.
\end{proof}

Karena model kanonik bagi \Log{K} menentukan~\Log{K}, kita langsung
memperoleh kelengkapan~\Log{K} sebagai suatu akibat:

\begin{cor}\ollabel{cor:Kcomplete}
  Logika modal dasar \Log{K} lengkap terhadap kelas semua model, yakni jika
  $\Entails !A$, maka $\Log{K} \Proves !A$.
\end{cor}

\begin{proof}
  Secara kontrapositif, jika $\Log{K} \Proves/ !A$, maka berdasarkan
  Penentuan $\mSat/{M^{\Log{K}}}{!A}$, sehingga $!A$ tidak valid.
\end{proof}

Dalam kasus umum kelengkapan suatu sistem $\Sigma$ terhadap suatu kelas
model, misalnya kelengkapan \Log{KTB4} terhadap kelas model refleksif,
simetris, dan transitif, penentuan saja tidak cukup. Kita juga harus
menunjukkan bahwa model kanonik bagi sistem $\Sigma$ merupakan anggota kelas
tersebut, dan hal ini tidak serta-merta mengikuti dari konstruksi model
kanonik---bahkan tidak selalu benar!

\end{document}
